\documentclass[runningheads]{llncs}

\usepackage[T1]{fontenc}
\usepackage[utf8]{inputenc}
\usepackage{booktabs}
\usepackage{amsmath}
\usepackage{graphicx}
\usepackage{url}
\usepackage[hidelinks]{hyperref}
\usepackage{xcolor}

\newcommand{\num}[1]{\mbox{#1}}

\begin{document}

\title{How Much of an Expected Goals Model Is in the Words?}
\subtitle{Recovering shot quality from free-text match commentary}

\author{Dheepak Karan}
\authorrunning{D. Karan}
\institute{Northeastern University, Boston MA, USA\\
\email{elumalaisanthakuma.d@northeastern.edu}}

\maketitle

\begin{abstract}
Expected goals (xG) models are built on shot coordinates, which are collected
by hand or by camera and are largely paywalled. Free-text match commentary is
neither: it is published for far more competitions than coordinates are, it
needs no key, and it describes each shot in a sentence. We ask how much of a
coordinate-based xG model survives if the only input is that sentence.

We parse \num{87,980} shots from ESPN commentary across six European
competitions into eighteen binary and ordinal fields, fit a logistic
regression on three Premier League seasons, and evaluate on a held-out fourth
(AUC \num{0.7709}). Joining \num{8,825} shots to StatsBomb's open Premier
League 2015/16 data, where both a commentary sentence and a hand-collected
coordinate xG exist for the same shot, the text model reaches AUC \num{0.7826}
against StatsBomb's \num{0.8118}: the words recover \textbf{90.6\%} of the
coordinate model's discrimination above chance. Per team-match the two agree at
$r=\num{0.869}$, which sits inside the band that commercial providers occupy
against each other. Handing a model the whole sentence instead of the extracted
fields is worth $-0.010$ AUC, so the extraction is not the bottleneck; the
missing 9.4\% is information the sentence never contained.

We then report a failure mode specific to text-derived features. Confirming
Robberechts and Davis, a decade-stale model loses only \num{0.0068} AUC and
transfers between leagues at no cost --- the \emph{ranking} is stable. The
\emph{level} is not: the shipped model over-predicts the held-out season by
11.9\%, which decomposes exactly into a 5.7\% fall in the league goal rate and
a 5.9\% shift in phrase mix. Two phrases carry the second half; unambiguous
phrases such as ``penalty'' do not move at all. Retraining cannot fix this, but
refitting a single intercept on shots already played reduces the error to
2.9\% without touching the ranking. A coordinate does not drift. A phrase does,
because a person writes it.

\keywords{Expected goals \and Football analytics \and Text as features \and
Concept drift \and Calibration}
\end{abstract}

\section{Introduction}

An expected goals model estimates the probability that a shot becomes a goal.
Every widely used version is built on where the shot was taken from: a
coordinate pair, a distance, an angle, and increasingly the positions of
everyone else in frame~\cite{mead2023xg,cavus2022explainable}. Those
coordinates are the expensive part. They are produced by human annotators or
tracking cameras, and outside a handful of open releases they are sold rather
than published.

Free-text commentary is the opposite. For each shot, a sentence like

\begin{quote}
\textit{``Attempt saved. Bukayo Saka (Arsenal) right footed shot from the
centre of the box is saved in the bottom left corner. Assisted by Martin
{\O}degaard with a through ball following a fast break.''}
\end{quote}

\noindent
is published live, without a key, for competitions that will never have open
coordinates. It plainly contains \emph{some} of what an xG model wants. The
question this paper answers is: how much?

That framing matters, because the answer is a number rather than a yes or no.
If commentary recovers a small fraction of a coordinate model, it is a
curiosity. If it recovers most of one, then xG becomes available in every
competition where somebody writes sentences about the football --- which is
most of them.

\subsection*{Contributions}

\begin{enumerate}
\item \textbf{A measurement.} On \num{8,825} shots for which both a commentary
sentence and a StatsBomb coordinate xG exist, a model reading only the sentence
recovers \textbf{90.6\%} of the coordinate model's discrimination above chance
(Section~\ref{sec:headline}).
\item \textbf{A scale for that number.} Agreement with StatsBomb per team-match
is $r=\num{0.869}$, comparable to how closely commercial providers agree with
each other (Section~\ref{sec:band}).
\item \textbf{An upper bound on extraction.} Giving models the full sentence
instead of eighteen extracted fields --- every 1--4 gram, and a sentence
embedding --- makes them \emph{worse} by \num{0.010} AUC. The remaining gap is
not a parsing failure (Section~\ref{sec:ceiling}).
\item \textbf{A drift result specific to text features.} The ranking is stable
over a decade, replicating~\cite{robberechts2020data} on a different feature
source; the calibration level is not, and we identify the two phrases
responsible, show a control phrase that does not move, and give a leak-free
in-season correction (Section~\ref{sec:drift}).
\item \textbf{A reproducible pipeline.} Collection, parsing, training,
evaluation and every figure in this paper rebuild from one
script~\cite{xgfromwords2026}.
\end{enumerate}

\section{Related Work}

\paragraph{Coordinate xG.}
Expected goals is two decades old and well surveyed. Recent work has focused on
explainability~\cite{cavus2022explainable}, player and position
correction~\cite{scholtes2024bayesxg}, and demonstrating downstream
value~\cite{mead2023xg}. All of it takes location as given.

\paragraph{Data requirements of xG.}
Robberechts and Davis~\cite{robberechts2020data} ask what data an xG model
actually needs. They find a basic logistic model converges at roughly
\num{6,000} shots, that training on less recent data costs ``only a negligible
performance hit'', that league-specific and cross-league models ``perform
equally'', and that data from different providers should not be mixed. They
argue calibration, not ranking, is the objective. This paper is in their debt
twice: our staleness and transfer results replicate theirs from a different
feature source, and our drift finding is only interesting because they
established the coordinate-side baseline.

\paragraph{Language and football.}
Text has been used to \emph{forecast} events from event
sequences~\cite{mendesneves2024forecasting}, and to \emph{explain} an xG model
by turning its coefficients into sentences~\cite{rahimian2025wordalisations}.
The latter is this paper's mirror image: they go from a model to words, we go
from words to a model. That both directions are being explored suggests the
interface between the two is worth studying; it does not make either one the
other.

\paragraph{Prior art on the input.}
A public hobby project~\cite{ballard2020moneyball} scrapes Premier League
website commentary and states directly that ``we can infer the shot position
from the text commentary'', deriving goal-conversion rates by shot category. To
our knowledge this is the first use of commentary text as a shot-quality input,
and it predates this work. It reports no fitted model, no held-out evaluation,
no comparison against a reference xG model, and no AUC, log loss or Brier
score. \textbf{The idea is not ours. The measurement is what we contribute.}

\paragraph{Provider agreement.}
Public comparisons of Opta, Understat, StatsBomb and Wyscout report match-level
xG correlations between \num{0.86} and \num{0.96} depending on the
pair~\cite{pythonfootball2024providers}. We use this as the scale against which
our own agreement with StatsBomb should be read.

\paragraph{Feature corroboration.}
American Soccer Analysis publish their xG feature
set~\cite{asa2015xg3}: distance, goal mouth available, headed, off a cross or
through ball, and pattern of play including corners, free kicks, fast breaks
and penalties, fitted with logistic regression. With the exception of the two
continuous geometric terms, this is almost exactly the field list a commentary
sentence yields --- independent evidence that the phrase-derived fields are the
right ones to want.

\section{Data}

\subsection{Commentary}

ESPN publish an Opta-derived English commentary feed through a public endpoint
requiring no key. We collected \num{3,569} matches across six competitions
(Table~\ref{tab:corpus}), yielding \num{87,980} parsed shots and \num{10,109}
goals.

\begin{table}[t]
\caption{The commentary corpus. Premier League spans 2015/16 and
2022/23--2025/26; the other five competitions cover 2022/23--2025/26.}
\label{tab:corpus}
\centering
\begin{tabular}{lrrrr}
\toprule
Competition & Matches & Shots & Goals & Goal rate \\
\midrule
Premier League & \num{1,892} & \num{47,431} & \num{5,489} & 11.6\% \\
La Liga        & 380 & \num{9,240} & \num{1,024} & 11.1\% \\
Serie A        & 380 & \num{9,065} & 922         & 10.2\% \\
Bundesliga     & 306 & \num{7,853} & 990         & 12.6\% \\
Ligue 1        & 305 & \num{7,391} & 863         & 11.7\% \\
Primeira Liga  & 306 & \num{7,000} & 821         & 11.7\% \\
\midrule
\textbf{Total} & \textbf{\num{3,569}} & \textbf{\num{87,980}} &
\textbf{\num{10,109}} & \textbf{11.5\%} \\
\bottomrule
\end{tabular}
\end{table}

\subsection{The reference model}

StatsBomb release Premier League 2015/16 openly~\cite{statsbomb2024open}, with
a true shot location, a freeze frame of every player, and their own xG for each
shot. Joining it to ESPN commentary for the same matches is the only way to put
words and coordinates on the same shots.

The join is deliberately conservative: a shot pairs only when the match, the
team and the minute all agree, and an unmatched shot is dropped rather than
guessed at. This yields \textbf{\num{8,825} shots across 373 matches}. As a
sanity check on the join itself, the two sources agree on whether the shot was
a goal for \textbf{99.71\%} of paired shots (26 disagreements). All evaluation
in Section~\ref{sec:headline} uses StatsBomb's goal flag as the label, so their
model is not penalised for annotation differences in ours.

\section{Method}

\subsection{The label is in the sentence}
\label{sec:leak}

Commentary states the outcome before it describes the chance. Every line begins
``Goal!'', ``Attempt missed'', ``Attempt saved'' or ``Attempt blocked'', and
often repeats it in the verb (``is saved in the bottom left corner'', ``hits
the left post''). That text \emph{is} the label. A model handed the raw
sentence scores a perfect AUC and has learned nothing.

Removing it is harder than it looks, and four separate leaks survived earlier
attempts:

\begin{enumerate}
\item Stripping only the opening clause left penalties at 100\% conversion.
\item A blacklist of outcome words missed the verbs; inspection of the fitted
coefficients found \texttt{goal} weighted at $+7.99$.
\item After moving to a \textbf{whitelist} --- keeping only spans matching a
fixed list of chance-describing patterns --- the phrase ``from a free kick''
still converted 60 of 60. Free kicks are consequently absent from the feature
set.
\item The \texttt{header} field was being set by the phrase ``headed
\emph{pass}'', mislabelling \num{1,636} shots and inflating their apparent
conversion from 10.0\% to 13.8\%. Fields describing the shot itself are now
read only from the text preceding ``assisted by''.
\end{enumerate}

The lesson generalises: with text features, a leak is a phrasing nobody thought
of, so the guard has to be behavioural rather than a word list. We train a
model on the stripped text and assert it \emph{cannot} separate goals too well
(ceiling AUC \num{0.85}, against \num{1.00} for the raw text), and we flag any
$n$-gram whose Wilson lower bound on conversion exceeds that of penalties ---
the best legitimate feature. This is the check that caught leak (3) and cleared
a genuine high-conversion phrase (``close range following fast break'',
26 of 32) that a fixed 80\% threshold had wrongly flagged.

\subsection{Features}

Eighteen fields are extracted by whitelist regex: six location zones
(\texttt{six\_yard}, \texttt{centre\_box}, \texttt{side\_box},
\texttt{outside\_box}, \texttt{long\_range}, \texttt{difficult\_ang}), three
body parts, five build-up descriptors (\texttt{from\_cross},
\texttt{from\_headed\_pass}, \texttt{from\_through}, \texttt{after\_corner},
\texttt{after\_break}), \texttt{assisted}, \texttt{penalty}, and the minute.

Table~\ref{tab:signal} shows the resulting spread. A twentyfold range in
conversion, recovered from English sentences alone, is the reason the rest of
the paper is possible.

\begin{table}[t]
\caption{Goal rate by parsed field, \num{37,929} Premier League shots
(2022/23--2025/26). Overall rate 11.8\%.}
\label{tab:signal}
\centering
\begin{tabular}{lrr@{\hskip 2em}lrr}
\toprule
Field & \% shots & Goals & Field & \% shots & Goals \\
\midrule
penalty          & 1.0  & 82.7\% & from cross      & 18.6 & 10.7\% \\
six yard box     & 5.0  & 39.6\% & assisted        & 75.9 & 10.7\% \\
after fast break & 2.7  & 35.5\% & header          & 18.2 & 10.5\% \\
through ball     & 3.5  & 25.7\% & difficult angle & 2.5  & 8.1\%  \\
after corner     & 9.8  & 16.4\% & side of box     & 18.5 & 5.6\%  \\
centre of box    & 36.2 & 14.4\% & outside the box & 30.6 & 4.1\%  \\
\bottomrule
\end{tabular}
\end{table}

\subsection{Model and protocol}

The shipped model is a standardised logistic regression. It is deliberately the
simplest thing that works: with eighteen mostly binary fields there is little
for a stronger learner to find, and Section~\ref{sec:ceiling} confirms this
empirically.

Two distinct fits are reported and should not be confused.

\begin{itemize}
\item \textbf{The shipped model}, trained on Premier League 2022/23--2024/25
(\num{28,735} shots) and evaluated on the held-out 2025/26 season
(\num{9,194} shots).
\item \textbf{The comparison model}, trained on Premier League
2022/23--2025/26 (\num{37,929} shots) and evaluated on 2015/16, which is
excluded from its training data. This is the fit used against StatsBomb.
\end{itemize}

All training is scoped to the Premier League, following
\cite{robberechts2020data} on not mixing sources; the other five competitions
are held out entirely as a transfer test (Section~\ref{sec:transfer}).

\section{Results}

\subsection{Held-out season}

On 2025/26, unseen during training, the model reaches
\textbf{AUC \num{0.7709}}, Brier \num{0.0840}, against a base rate of 11.3\%.

\subsection{Against a coordinate model}
\label{sec:headline}

Table~\ref{tab:headline} is the central result. Both models predict the same
\num{8,825} shots, labelled by StatsBomb.

\begin{table}[t]
\caption{Words against coordinates, \num{8,825} shots, Premier League 2015/16.}
\label{tab:headline}
\centering
\begin{tabular}{llrrr}
\toprule
Model & Input & AUC & Log loss & Brier \\
\midrule
Ours & the commentary sentence & \num{0.7826} & \num{0.2711} & \num{0.0764} \\
StatsBomb & coordinates, freeze frames, 16 player positions
& \num{0.8118} & \num{0.2555} & \num{0.0715} \\
\bottomrule
\end{tabular}
\end{table}

\begin{equation}
\frac{0.7826 - 0.5}{0.8118 - 0.5} = \mathbf{0.906}
\end{equation}

\noindent
\textbf{The words recover 90.6\% of the coordinate model's discrimination above
chance.} The mean predictions agree to within \num{0.001} (\num{0.0972} against
\num{0.0982}), so whatever the words miss, they do not miss it in a biased
direction. Per-shot correlation is \num{0.735} (Spearman \num{0.628}), mean
absolute difference \num{0.052}.

\subsection{Is 90.6\% close? A scale}
\label{sec:band}

A recovery figure means little without something to compare it to. The natural
comparison is how closely the commercial providers agree with \emph{each
other}~\cite{pythonfootball2024providers}, and Table~\ref{tab:band} places our
model on the same axis.

\begin{table}[t]
\caption{Match-level xG agreement. Provider figures
from~\cite{pythonfootball2024providers} (five leagues, five seasons,
$\approx$\num{4,290} matches); ours computed on 746 team-matches of the join.}
\label{tab:band}
\centering
\begin{tabular}{lc}
\toprule
Pair & Correlation \\
\midrule
Opta $\times$ Understat & \num{0.96} \\
Opta $\times$ StatsBomb & \num{0.92}--\num{0.93} \\
StatsBomb $\times$ Understat & \num{0.92}--\num{0.93} \\
Wyscout $\times$ the others & \num{0.86}--\num{0.88} \\
\midrule
\textbf{Ours (commentary text) $\times$ StatsBomb} & \textbf{\num{0.869}} \\
\bottomrule
\end{tabular}
\end{table}

Mean absolute difference is \num{0.249} xG per team-match (median \num{0.198},
max \num{1.54}); mean xG \num{1.150} against StatsBomb's \num{1.161} and
\num{1.192} actual goals.

\textbf{A model that reads only English sentences agrees with StatsBomb about
as closely as a commercial provider does.} Two caveats attach. First, the
aggregation levels are close but not identical: theirs is per match across five
leagues and five seasons, ours is per team-match on one league and one season.
Second, a widely repeated ``$\approx$1 xG mean absolute difference between
providers, maximum \num{3.88}'' figure is deliberately \emph{not} used here.
Naming a single club alongside \num{3.88} reads as a season aggregate, which
would make that figure far tighter per match than our \num{0.249} rather than
looser --- and the source paywalls, so we cannot confirm the unit.

\subsection{Is the extraction the bottleneck? No}
\label{sec:ceiling}

The 9.4\% gap has two possible homes: the regexes catch only phrasings someone
thought of, or the sentence never contained the information. This is decided by
handing models the whole sentence and seeing whether they beat the eighteen
fields pulled out of it (Table~\ref{tab:ceiling}).

\begin{table}[t]
\caption{Reading more of the sentence, held-out 2025/26 season.}
\label{tab:ceiling}
\centering
\begin{tabular}{llr}
\toprule
Model & Sees & AUC \\
\midrule
\textbf{18 regex fields, logistic (shipped)} & 18 fields & \textbf{\num{0.7709}} \\
18 regex fields, gradient-boosted trees & 18 fields & \num{0.7695} \\
every 1--4 gram, TF-IDF & all words & \num{0.7612} \\
sentence embedding (MiniLM) & all words & \num{0.7584} \\
embedding $+$ regex fields & both & \num{0.7589} \\
\bottomrule
\end{tabular}
\end{table}

Every model given the full text does \emph{worse}, including a semantic
embedding; adding the embedding to the fields makes those worse too. The rest
of the sentence is player names, team names and filler. \textbf{Reading more is
worth $-0.010$ AUC}, so no better reader --- an LLM extractor included ---
closes the gap. The missing 9.4\% is information the sentence never contained:
the exact distance inside a zone, the angle, and how many defenders stood in
the way.

\subsection{Transfer}
\label{sec:transfer}

The claim that xG becomes available wherever commentary exists is only real if
the model travels. The Premier League model was pointed at five other
competitions cold --- nothing retrained, nothing tuned (Table~\ref{tab:transfer}).

\begin{table}[t]
\caption{One model, six competitions. Nothing retrained.}
\label{tab:transfer}
\centering
\begin{tabular}{lrrrr}
\toprule
Competition & Shots & Goal rate & AUC & Mean xG \\
\midrule
Premier League \emph{(trained on)} & \num{9,194} & 11.3\% & \num{0.7709} & \num{0.127} \\
La Liga        & \num{9,240} & 11.1\% & \num{0.7730} & \num{0.119} \\
Ligue 1        & \num{7,391} & 11.7\% & \num{0.7807} & \num{0.122} \\
Bundesliga     & \num{7,853} & 12.6\% & \num{0.7795} & \num{0.125} \\
Serie A        & \num{9,065} & 10.2\% & \num{0.7702} & \num{0.115} \\
Primeira Liga  & \num{7,000} & 11.7\% & \num{0.7871} & \num{0.122} \\
\bottomrule
\end{tabular}
\end{table}

Mean AUC abroad is \num{0.7781} against \num{0.7709} at home: a transfer cost
of $-0.0072$, which is to say a small gain. Calibration holds, with predictions
of \num{0.115}--\num{0.125} against observed rates of 10.2--12.6\% and the
ordering preserved. This replicates~\cite{robberechts2020data} on a feature
source they did not consider, and it works for a reason specific to this
setting: Opta build the sentence the same way in every competition.

\section{The Ranking Holds, the Level Does Not}
\label{sec:drift}

\subsection{The symptom}

On the held-out season the model predicts \num{1.54} goals per team-match
against \num{1.38} actual --- an \textbf{11.9\% over-estimate}. It decomposes
exactly:

\begin{center}
\begin{tabular}{lcr}
\toprule
& & Factor \\
\midrule
league goal rate, 2022/23 $\rightarrow$ 2025/26
  & $0.1199 \rightarrow 0.1134$ & $\times\,1.057$ \\
mean prediction vs.\ its own training rate
  & $0.1199 \rightarrow 0.1270$ & $\times\,1.059$ \\
\midrule
\textbf{observed over-estimate} & & $\mathbf{1.119}$ \\
\bottomrule
\end{tabular}
\end{center}

The first factor is the league converting slightly less often. The second is
stranger: a logistic regression's mean prediction on its training set equals the
training base rate by construction, so a mean of \num{0.1270} on 2025/26 means
the \emph{parsed features} of those shots look better than the shots the model
was fitted on --- while producing fewer goals.

\subsection{Two phrases, and a control that does not move}

\begin{table}[t]
\caption{Phrase drift in Premier League commentary. Two phrases move; two
controls do not.}
\label{tab:drift}
\centering
\begin{tabular}{lrrrrr}
\toprule
& \multicolumn{2}{c}{2015/16} & \multicolumn{2}{c}{2025/26} & \\
\cmidrule(lr){2-3}\cmidrule(lr){4-5}
Phrase & \% shots & conv. & \% shots & conv. & \\
\midrule
``six yard box''        & 3.22 & 43.5\% & 5.96 & 35.8\% & \textbf{drifts} \\
``following a fast break'' & 0.73 & 49.3\% & 3.30 & 28.4\% & \textbf{drifts} \\
\midrule
``centre of the box''   & 32.84 & 14.2\% & 36.44 & 13.2\% & stable \\
``penalty''             & 0.94 & 83.1\% & 1.01 & 82.8\% & stable \\
\bottomrule
\end{tabular}
\end{table}

Table~\ref{tab:drift} localises the second factor. ``Following a fast break''
became \textbf{4.5 times more common at little over half the conversion} in a
decade. That is not a change in football. It is a change in how loosely the
phrase is written, and the model --- which knows nothing except the words ---
reads a better shot than the one that was taken.

The controls matter as much as the finding. ``Penalty'' is unambiguous, and it
does not move: 0.94\% of shots at 83.1\% in 2015/16, 1.01\% at 82.8\% in
2025/26. ``Centre of the box'' is nearly as stable. \textbf{The drift is not a
general artefact of the parser or the corpus; it is specific to phrases that
require a judgement call from whoever is writing.}

Scored with the shipped model, the sign of the error flips across the decade:
$-2.7\%$ on 2015/16, $+11.9\%$ on 2025/26. The model was fitted in between.

\subsection{Retraining does not help}

\begin{center}
\begin{tabular}{lr}
\toprule
Training window & Over-estimate \\
\midrule
all three seasons (shipped) & $+11.9\%$ \\
last two seasons & $+11.9\%$ \\
last season only & $+9.1\%$ \\
recency weighted, half-life one season & $+9.9\%$ \\
\bottomrule
\end{tabular}
\end{center}

The best of these buys 2.8 of the 11.9 points and pays two thirds of the
training set for it. The reason is not that the data is old: \textbf{next
season's conversion rate is not in this season's data at any weighting.}

\subsection{Waiting does}

Once a few hundred shots of the new season have been played, their realised
rate is observable, and one number added to the intercept pulls the level back.
We walk forward through 2025/26 in blocks of 500 shots, refitting the shift by
bisection on \emph{only} the shots played before each block:

\begin{center}
\begin{tabular}{lrrrr}
\toprule
& mean xG & actual & over & log loss \\
\midrule
uncorrected  & \num{0.1267} & \num{0.1132} & $+12.0\%$ & \num{0.29523} \\
recalibrated & \num{0.1165} & \num{0.1132} & $+2.9\%$  & \num{0.29443} \\
\bottomrule
\end{tabular}
\end{center}

\num{8,694} shots evaluated after a 500-shot warm-up --- roughly twenty
matches --- during which the model runs uncorrected because the realised rate
is still too noisy to correct towards. An intercept shift is monotone, so AUC
is unchanged by construction, and our test suite asserts it: the guard is
against improving calibration by reaching for the coefficients, which would
trade away the ranking to buy a better mean.

\subsection{Why this is not a contradiction of prior work}

Robberechts and Davis~\cite{robberechts2020data} put calibration at the centre
of xG evaluation and found staleness negligible. We reproduce their ranking
result exactly --- a decade-stale model costs \num{0.0068} AUC --- and observe
the level moving 10\% over the same span. Both are true, and together they say
something neither says alone:

\begin{quote}
\textbf{A coordinate does not drift. A phrase does, because a person writes
it.} Event-data xG is stable in precisely the sense --- calibration --- where a
text-derived model is not, and the instability is confined to phrases requiring
human judgement.
\end{quote}

This is, to our knowledge, the first characterisation of concept drift in
text-derived sports features, and it has a practical consequence: any system
built on commentary phrasing needs an in-season level correction, and does not
need retraining.

\section{Limitations}

We state these in full rather than in passing.

\begin{itemize}
\item \textbf{One reference model, one season.} The headline compares against
StatsBomb alone, on Premier League 2015/16. A second provider, or a second
season, would strengthen it considerably. StatsBomb's open releases are the
constraint.

\item \textbf{Free kicks are absent.} The phrase ``from a free kick'' leaked
(60 of 60 conversions), and rather than repair it we removed the feature. A
model that handled free kicks properly would likely do better.

\item \textbf{Matches with no commentary are dropped silently.} A small number
of matches arrive with an empty feed and contribute nothing. The rate is low
but it is not zero and it is not currently guarded.

\item \textbf{Latency is unmeasured.} We claim commentary arrives fast enough
for live use and have no measurement to support it. A scheduled matchday
watcher exists but the cloud scheduler that triggers it fired 2 of 12 slots on
the day tested, and missed the kickoff.

\item \textbf{The novelty claim rests on an imperfect search.} Our literature
search of arXiv, OpenAlex and Crossref returned 5 of 6 recall probes and not
always the same five. A follow-up blog and grey-literature search --- which
found~\cite{ballard2020moneyball} --- returned 3 of 4 probes by name, and
\texttt{reddit.com} refuses our crawler entirely, which is exactly where an
unpublished version of this would be discussed. Both searches are recorded
query-by-query in the repository. We have not searched paywalled indexes,
thesis repositories, or non-English sources.

\item \textbf{The provider-agreement comparison is close but not exact}, as
noted in Section~\ref{sec:band}.
\end{itemize}

\section{Conclusion}

A sentence of English recovers 90.6\% of a tracking-data expected goals model's
discrimination above chance, and agrees with it per team-match about as closely
as one commercial provider agrees with another. The gap is not a parsing
failure --- richer readers of the same sentence do worse --- but information
the sentence never contained. The model transfers across six competitions at no
cost, because the sentence is built the same way in each.

The interesting failure is not in the ranking but in the level. Text features
drift where coordinate features do not, because a phrase is written by a person
whose habits change, and we localise that drift to phrases requiring judgement
while unambiguous phrases stay fixed. One intercept, refitted in season on
shots already played, corrects it.

For competitions that will never have open coordinates --- which is most of
them --- this is the difference between having an xG model and not having one.

\begin{credits}
\subsubsection{\ackname}
This work uses StatsBomb open data and a public ESPN endpoint. No funding was
received.

\subsubsection{\discintname}
The author declares no competing interests.
\end{credits}

\section*{Reproducibility}

Every number in this paper is produced by \texttt{rebuild.sh}
in~\cite{xgfromwords2026}, which runs the pipeline in dependency order from raw
collection to final figures. The repository additionally contains a test suite
of 50 tests, including behavioural leak guards, a walk-forward check that the
recalibration uses no future shots, an assertion that the shift leaves AUC
unchanged, and a dependency graph that fails if any artefact predates its
inputs. The literature and blog searches are recorded query-by-query with their
recall probes.

\bibliographystyle{splncs04}
\bibliography{refs}

\end{document}
